\scanpage{81}
% archive foliation: 80
Ainsi
\[X\longrightarrow\Gamma^*(X)=(\Gamma^i(X))_{i\geq0}\]
est un foncteur (pas \inserted{\unsure{néc.}} additif !)
\[\unsure{\mathcal P|P}_1\longrightarrow\operatorname{Grad}(\unsure{\mathcal P|P}_1)\]
qui est en fait un foncteur ``transformant sommes en $\otimes$''. Comme tout objet de $\unsure{\mathcal P|P}_1$ est un \inserted{\unsure{objet}} groupe comm. pour l'opération $\oplus$, il s'ensuit que dans chaque $\Gamma^*(X)$ il y a une structure $\otimes$-algèbre \inserted{graduée} comm., fonctorielle en $X$.
\margin{Pour $\unsure{C=|\Gamma}$ \unsure{voir que} \unsure{ces assertions} \unsure{résultent des} \deleted{\unsure{bigradués}} \unsure{avec en fait} $2$ \unsure{structures}.}
\deleted{En particulier}

On en définit \inserted{dans $\unsure{\mathcal P|P}_1$} un hom, fonctoriel en $X$
\[\text{(\unsure{\alpha})}\qquad\Gamma^p(\Gamma^q(X))\longleftarrow\Gamma^{pq}(X)\]
ou ce qui revient au m., un hom de degré $pq$
\[X\longrightarrow\Gamma^p(\Gamma^q(X)).\]
On prend le composé
\diagram{\begin{tikzcd}[column sep=large]
X\arrow[r,"\unsure{e}^X_q","{\substack{\text{hom univ.}\\\text{de degré }q}}"']&\Gamma^q(X)\arrow[r,"\unsure{e}^{\Gamma^q(X)}_p","{\substack{\text{hom univ.}\\\text{de degré }p}}"']&\Gamma^p(\Gamma^q(X))
\end{tikzcd}}
On a la condition ``d'associativité''
\diagram{\begin{tikzcd}[column sep=huge,row sep=large]
\text{(A)}&\Gamma^{pqr}(X)\arrow[r,"{\alpha^X_{\unsure{p,qr}}}"]\arrow[d,"{\alpha^X_{pq,r}}"']&\Gamma^p(\Gamma^{qr}(X))\arrow[d,"{\Gamma^p(\alpha^{\unsure{\Gamma^r(X)}}_{q,r})}"]\\
&\Gamma^{pq}(\Gamma^r(X))\arrow[r,"{\deleted{\unsure{\Gamma}}\alpha^{\Gamma^r(X)}_{p,q}}"']\arrow[ur,phantom,"\text{comm.}" description]&\Gamma^p(\Gamma^q(\Gamma^r(X)))
\end{tikzcd}}
et d'\uline{unité} : l'\unsure{isom.} $\unsure{\alpha}$ est l'identité \unsure{car pour} $\Gamma^1X=X$, $\Gamma^0X=1$

